
\documentclass{article}
\usepackage{colortbl}
\usepackage{makecell}
\usepackage{multirow}
\usepackage{supertabular}

\begin{document}

\newcounter{utterance}

\centering \large Interaction Transcript for game `cladder', experiment `full\_v1.5\_default', episode 2646 with Qwen3.5{-}4B.
\vspace{24pt}

{ \footnotesize  \setcounter{utterance}{1}
\setlength{\tabcolsep}{0pt}
\begin{supertabular}{c@{$\;$}|p{.15\linewidth}@{}p{.15\linewidth}p{.15\linewidth}p{.15\linewidth}p{.15\linewidth}p{.15\linewidth}}
    \# & \multicolumn{2}{c}{Player} && \multicolumn{2}{c}{Game Master} \\
    \hline

    \theutterance \stepcounter{utterance}  
    & & & \multicolumn{4}{p{0.6\linewidth}}{
        \cellcolor[rgb]{0.9,0.9,0.9}{
            \makecell[{{p{\linewidth}}}]{
                \texttt{\tiny{[P1$\langle$GM]}}
                \texttt{You are an expert in causal inference. The following question is not a typical commonsense query, but rather a meticulously designed question created by a professor specializing in causal inference, intended to assess the students' mastery of the course content.} \\
\texttt{Imagine a self{-}contained, hypothetical world with only the following conditions, and without any unmentioned factors or causal relationships: Gender has a direct effect on smoking and lung cancer. Smoking has a direct effect on tar deposit. Tar deposit has a direct effect on lung cancer. Gender is unobserved. For nonsmokers and with no tar deposit, the probability of lung cancer is 79\%. For nonsmokers and with high tar deposit, the probability of lung cancer is 49\%. For smokers and with no tar deposit, the probability of lung cancer is 71\%. For smokers and with high tar deposit, the probability of lung cancer is 40\%. For nonsmokers, the probability of high tar deposit is 79\%. For smokers, the probability of high tar deposit is 40\%. For smokers, would it be more likely to see lung cancer if the person had been a nonsmoker?} \\
\texttt{Start your answer with 'Yes' or 'No', followed by additional reasoning or evidence to support your explanation.} \\
            }
        }
    }
     \\ \\

    \theutterance \stepcounter{utterance}  
    & \multicolumn{4}{p{0.6\linewidth}}{
        \cellcolor[rgb]{0.561,0.737,0.561}{
            \makecell[{{p{\linewidth}}}]{
                \texttt{\tiny{[P1$\rangle$GM]}}
                \texttt{Yes} \\
\\ 
\texttt{**Reasoning:**} \\
\\ 
\texttt{To determine whether a smoker would be more likely to see lung cancer if they had been a nonsmoker, we must calculate the counterfactual probability of lung cancer for a smoker under the condition that they were a nonsmoker, while holding all other causal factors constant according to the rules of this hypothetical world.} \\
\\ 
\texttt{1.  **Identify the Causal Pathways:**} \\
\texttt{    *   Gender \$\textbackslash{}rightarrow\$ Smoking} \\
\texttt{    *   Gender \$\textbackslash{}rightarrow\$ Tar Deposit (implied by the direct effect mentioned in the prompt alongside smoking, though the prompt explicitly states "Gender has a direct effect on smoking and lung cancer" and "Smoking has a direct effect on tar deposit" and "Tar deposit has a direct effect on lung cancer". The prompt also gives probabilities for tar deposits based on smoking status, implying Gender influences the probability of being a smoker, which then influences tar deposits. However, the critical constraint is the conditional probability of lung cancer given the state of the person).} \\
\texttt{    *   Smoking \$\textbackslash{}rightarrow\$ Tar Deposit} \\
\texttt{    *   Tar Deposit \$\textbackslash{}rightarrow\$ Lung Cancer} \\
\texttt{    *   Smoking \$\textbackslash{}rightarrow\$ Lung Cancer (directly, as stated: "Gender has a direct effect on smoking and lung cancer" implies a direct link from Gender to Lung Cancer, but the prompt also gives specific probabilities for "Smokers" and "Nonsmokers" regarding lung cancer outcomes, which likely aggregates the direct effect of gender on lung cancer plus the indirect effects via smoking and tar. Let's re{-}read carefully: "Gender has a direct effect on smoking and lung cancer." This usually implies a backdoor path or a direct arrow. However, the data provided is conditional on "Smokers" and "Nonsmokers". Let's look at the specific probabilities given for the final outcome).} \\
\\ 
\texttt{    Actually, let's look at the specific conditional probabilities provided for the outcome "Lung Cancer":} \\
\texttt{    *   \$P(\textbackslash{}text\{Lung Cancer\} | \textbackslash{}text\{Nonsmoker, No Tar\}) = 0.79\$} \\
\texttt{    *   \$P(\textbackslash{}text\{Lung Cancer\} | \textbackslash{}text\{Nonsmoker, High Tar\}) = 0.49\$} \\
\texttt{    *   \$P(\textbackslash{}text\{Lung Cancer\} | \textbackslash{}text\{Smoker, No Tar\}) = 0.71\$} \\
\texttt{    *   \$P(\textbackslash{}text\{Lung Cancer\} | \textbackslash{}text\{Smoker, High Tar\}) = 0.40\$} \\
\\ 
\texttt{    The question asks: "For smokers, would it be more likely to see lung cancer if the person had been a nonsmoker?"} \\
\texttt{    This is a counterfactual query: \$P(\textbackslash{}text\{Lung Cancer\} | \textbackslash{}text\{Smoker\}, \textbackslash{}text\{do\}(\textbackslash{}text\{Nonsmoker\}))\$.} \\
\\ 
\texttt{2.  **Apply the Do{-}Calculus / Counterfactual Logic} \\
            }
        }
    }
    & & \\ \\

    \theutterance \stepcounter{utterance}  
    & & & \multicolumn{2}{p{0.3\linewidth}}{
        \cellcolor[rgb]{0.9,0.9,0.9}{
            \makecell[{{p{\linewidth}}}]{
                \texttt{\tiny{[GM$|$GM]}}
                \texttt{Yes} \\
            }
        }
    }
    & & \\ \\

    \theutterance \stepcounter{utterance}  
    & & & \multicolumn{2}{p{0.3\linewidth}}{
        \cellcolor[rgb]{0.9,0.9,0.9}{
            \makecell[{{p{\linewidth}}}]{
                \texttt{\tiny{[GM$|$GM]}}
                \texttt{no} \\
            }
        }
    }
    & & \\ \\

    \theutterance \stepcounter{utterance}  
    & & & \multicolumn{2}{p{0.3\linewidth}}{
        \cellcolor[rgb]{0.9,0.9,0.9}{
            \makecell[{{p{\linewidth}}}]{
                \texttt{\tiny{[GM$|$GM]}}
                \texttt{game\_result = LOSE} \\
            }
        }
    }
    & & \\ \\

\end{supertabular}
}

\end{document}
